%% ============================================================
%%  Abstract (English)
%% ============================================================
\chapter*{Abstract}
\addcontentsline{toc}{chapter}{Abstract}
\thispagestyle{plain}

\begin{english}
Quantitative trading on the Casablanca Stock Exchange (MASI) faces two structural
challenges: rapidly identifying the best market opportunities in an information-dense
environment, and rigorously validating strategies against overfitting. This final-year
engineering project, conducted at the \textbf{BMCE Capital Trading Desk}, addresses
both challenges through the design and development of a full quantitative platform
organized as a four-station production line.

\medskip

The platform implements two complementary result engines. The first, the
\textit{Signal Engine}, evaluates 20 families of technical indicators across 93 MASI
tickers through a seven-layer pipeline (Layers A→G) based on rolling out-of-sample
(OOS) evaluation, a multi-component robustness score, and redundancy reduction via
correlation clustering. The second, the \textit{WFO Backtest Engine}, applies Pardo's
(2008) Walk-Forward Analysis methodology with the PROM objective function, producing
a robustness verdict through Walk-Forward Efficiency (WFE).

\medskip

Both engines feed into a \textbf{unified dashboard} enabling traders to move from
market reading to investment decision — direction, conviction, actionable levels —
in under 30 seconds per ticker.

\medskip

Technically, the platform is built on a full-stack architecture (Next.js, FastAPI,
PostgreSQL, MinIO, RQ) with a Numba JIT-optimized quantitative core achieving a 38×
speedup on the backtesting kernel.

\bigskip
\noindent\textbf{Keywords:}
quantitative backtesting, walk-forward analysis, signal engine, technical indicators,
sub-period validation, MASI, algorithmic trading, overfitting, PROM, Numba.
\end{english}
